\documentclass[11pt]{article}
\usepackage[margin=1in]{geometry}
\usepackage[T1]{fontenc}
\usepackage{lmodern}
\usepackage{amsmath,amssymb,mathtools,booktabs,longtable,array,hyperref}
\hypersetup{colorlinks=true,linkcolor=blue,urlcolor=blue,citecolor=blue}

\newcommand{\Lmu}{L_\mu}
\newcommand{\Span}{\operatorname{span}}
\newcommand{\dd}{\mathrm{d}}

\title{Step-by-Step Derivation of the ORBIT Order-6 Equivalence}
\author{ORBIT}
\date{2026-03-23}

\begin{document}
\maketitle
\tableofcontents

\section{Purpose}

This note rewrites the stored ORBIT order-6 comparison as a paper-style derivation. The point is not to quote the code, but to repeat the same algorithm mathematically:
\begin{enumerate}
\item reduce both Lagrangians sector by sector to the same IBP / field-redefinition quotient basis,
\item canonically normalize the quadratic kinetic term,
\item solve the reduced-coordinate matching equations,
\item isolate the only leftover sector $\{2,4\}$,
\item show explicitly that this leftover is a quadratic higher-derivative field-redefinition image.
\end{enumerate}

The companion raw audit report
\texttt{orbit\_fixed\_eh\_equivalence\_audit\_report.pdf}
contains the verbatim stored expressions, source excerpts, and exact zero checks used below.

\section{Setup}

The comparison is between:
\begin{itemize}
\item the Matchete-derived EFT, translated to xAct and then reduced by ORBIT,
\item the reference sign-flipped Einstein--Hilbert Lagrangian with cosmological term,
\[
\mathcal{L}_{\mathrm{ref}}
=
\Lambda_{\mathrm{ref}}\sqrt{-g}
-\frac{2}{\kappa_{\mathrm{ref}}^2}\sqrt{-g}\,R,
\]
expanded around mostly-minus flat space with
\[
g_{\mu\nu}=\eta_{\mu\nu}+\kappa_{\mathrm{ref}} h_{\mu\nu}.
\]
\end{itemize}

Throughout, define
\[
\Lmu := \log\!\frac{\bar\mu^2}{M^2}.
\]

ORBIT decomposes the scalar operator space into sectors
\[
V_{n_h,N_\partial},
\]
where $n_h$ is the graviton number and $N_\partial$ is the derivative number. In each sector one chooses a reduced basis
\[
B_{n_h,N_\partial}=\{b_1,\dots,b_r\},
\]
after quotienting by:
\begin{itemize}
\item integration by parts in every sector,
\item the Fierz--Pauli field-redefinition image in interaction sectors.
\end{itemize}

If the reduced coordinates of the input and reference are
\[
v^{\mathrm{in}}_{n_h,N_\partial},
\qquad
v^{\mathrm{ref}}_{n_h,N_\partial},
\]
then the reduced difference is
\[
\Delta v_{n_h,N_\partial}
:=
v^{\mathrm{in}}_{n_h,N_\partial}
-v^{\mathrm{ref}}_{n_h,N_\partial}.
\]
Equality in the implemented quotient means $\Delta v_{n_h,N_\partial}=0$ in every supported sector.

\section{Step 1: Canonical Normalization}

The reduced quadratic two-derivative coordinates of the input are proportional to the Fierz--Pauli coordinates. ORBIT stores
\[
v^{\mathrm{FP}}_{2,2}
=
\left(\frac12,-1,1,-\frac12\right),
\]
and the input coordinates simplify exactly to
\[
v^{\mathrm{in}}_{2,2}
=
Z_h\,v^{\mathrm{FP}}_{2,2},
\qquad
Z_h=\frac{M^2\kappa^2(1+\Lmu)}{24}.
\]

The final comparison uses target kinetic sign $-1$, so the canonically normalized field is defined by
\[
h_{\mu\nu}=\beta\,\widehat h_{\mu\nu},
\qquad
Z_h\,\beta^2=-1.
\]
Therefore
\[
\beta
=
\sqrt{-\frac{1}{Z_h}}
=
\sqrt{-\frac{24}{M^2\kappa^2(1+\Lmu)}}.
\]

It is convenient to define
\[
s:=\sqrt{-\frac{1}{M^2\kappa^2(1+\Lmu)}},
\]
so that
\[
\beta=2\sqrt6\,s.
\]

\section{Step 2: Reduced Matching Equations}

The full stored run produces 29 reduced-coordinate equations. For the actual derivation, only a small independent subset is needed.

\subsection{The equations that determine the branch}

From the reduced $\{6,0\}$ sector one obtains an equation with no reference parameters:
\begin{equation}
-864-576\Lmu=0.
\end{equation}
Hence
\begin{equation}
\boxed{\Lmu=-\frac32.}
\end{equation}

From the reduced $\{1,0\}$ sector:
\begin{equation}
-4\kappa_{\mathrm{ref}}\Lambda_{\mathrm{ref}}
\sqrt6\,M^4\kappa\,s\,(3+2\Lmu)=0.
\end{equation}
Using $\Lmu=-3/2$, the second term vanishes, so
\begin{equation}
\kappa_{\mathrm{ref}}\Lambda_{\mathrm{ref}}=0.
\end{equation}
The reference graviton coupling is nonzero on the branch of interest, therefore
\begin{equation}
\boxed{\Lambda_{\mathrm{ref}}=0.}
\end{equation}

From the reduced $\{4,2\}$ sector:
\begin{equation}
-24-\kappa_{\mathrm{ref}}^2 M^2(1+\Lmu)=0.
\end{equation}
Substituting $\Lmu=-3/2$ gives
\[
-24+\frac12 \kappa_{\mathrm{ref}}^2 M^2=0,
\]
hence
\begin{equation}
\boxed{\kappa_{\mathrm{ref}}^2=\frac{48}{M^2}.}
\end{equation}

From the reduced $\{3,2\}$ sector:
\begin{equation}
\kappa_{\mathrm{ref}}-2\sqrt6\,\kappa\,s=0.
\end{equation}
Since $\beta=2\sqrt6\,s$, this is
\[
\kappa_{\mathrm{ref}}=\kappa\beta.
\]
On the positive branch used in the stored verification,
\[
s=\sqrt{-\frac{1}{M^2\kappa^2(1+\Lmu)}}\Bigg|_{\Lmu=-3/2}
=
\frac{\sqrt2}{M\kappa},
\]
so
\begin{equation}
\boxed{
\beta=\frac{4\sqrt3}{M\kappa},
\qquad
\kappa_{\mathrm{ref}}=\frac{4\sqrt3}{M}.}
\end{equation}

Equivalently,
\begin{equation}
\boxed{M^2=e^{3/2}\bar\mu^2,}
\end{equation}
since $\Lmu=-3/2$ means $\bar\mu^2/M^2=e^{-3/2}$.

\subsection{Consistency of the remaining supported-sector equations}

Once
\[
\Lmu=-\frac32,\qquad
\Lambda_{\mathrm{ref}}=0,\qquad
\kappa_{\mathrm{ref}}=\frac{4\sqrt3}{M},
\]
are imposed, all remaining supported-sector equations vanish identically.

The reason is structural:
\begin{itemize}
\item the zero-derivative sectors $\{1,0\},\{2,0\},\{3,0\},\{4,0\},\{5,0\},\{6,0\}$ reduce to combinations of the factors $(3+2\Lmu)$ and $\Lambda_{\mathrm{ref}}$,
\item the derivative interaction sectors $\{3,2\},\{4,2\}$ reduce to combinations of
\[
\kappa_{\mathrm{ref}}-\kappa\beta,
\qquad
\kappa_{\mathrm{ref}}^2-\frac{48}{M^2},
\]
\item the quadratic kinetic sector $\{2,2\}$ was already made identical by canonical normalization.
\end{itemize}

\section{Step 3: Sector-by-Sector Vanishing After the Branch Is Imposed}

The stored exact zero checks in the companion audit report confirm that every supported sector vanishes after the substitutions
\[
\Lmu\to -\frac32,
\qquad
\Lambda_{\mathrm{ref}}\to 0,
\qquad
\kappa_{\mathrm{ref}}\to \frac{4\sqrt3}{M}.
\]

\begin{longtable}{@{}>{\raggedright\arraybackslash}p{1.3in}>{\raggedright\arraybackslash}p{4.8in}@{}}
\toprule
Sector & Reason the difference vanishes \\
\midrule
\endhead
$\{1,0\}$ & proportional to $-4\kappa_{\mathrm{ref}}\Lambda_{\mathrm{ref}}+\sqrt6 M^4\kappa s(3+2\Lmu)$ \\[3pt]
$\{2,0\}$ & proportional to $\kappa_{\mathrm{ref}}^2\Lambda_{\mathrm{ref}}(1+\Lmu)+3M^2(3+2\Lmu)$ \\[3pt]
$\{2,2\}$ & exactly equal after canonical normalization \\[3pt]
$\{3,0\}$ & proportional to $-\kappa_{\mathrm{ref}}^3\Lambda_{\mathrm{ref}}(1+\Lmu)+12\sqrt6 M^2\kappa s(3+2\Lmu)$ \\[3pt]
$\{3,2\}$ & proportional to $\kappa_{\mathrm{ref}}-\kappa\beta$ \\[3pt]
$\{4,0\}$ & proportional to $\kappa_{\mathrm{ref}}^4\Lambda_{\mathrm{ref}}(1+\Lmu)^2-72(3+2\Lmu)$ \\[3pt]
$\{4,2\}$ & proportional to $\kappa_{\mathrm{ref}}^2M^2(1+\Lmu)+24$ \\[3pt]
$\{5,0\}$ & every coordinate contains $(3+2\Lmu)$ and/or $\Lambda_{\mathrm{ref}}$ \\[3pt]
$\{6,0\}$ & every coordinate contains $(3+2\Lmu)$ and/or $\Lambda_{\mathrm{ref}}$ \\[3pt]
\bottomrule
\end{longtable}

Thus, after canonical normalization and the solved branch, every supported sector agrees with pure Einstein--Hilbert.

\section{Step 4: The Only Leftover Sector \texorpdfstring{$\{2,4\}$}{\{2,4\}}}

Before the final completion, the only unsupported difference is the quadratic four-derivative sector. ORBIT reduces it to a five-dimensional coordinate vector
\begin{equation}
v_{2,4}
=
\frac{1}{M^2}
\left(
\frac{9}{10},
-\frac{9}{20},
-\frac35,
\frac{3}{10},
-\frac{3}{20}
\right).
\end{equation}

To decide whether this is a genuine curvature-squared invariant or only a quadratic higher-derivative redefinition, ORBIT compares it against:
\begin{itemize}
\item the linearized curvature-squared vectors,
\item the quadratic higher-derivative redefinition-image vectors.
\end{itemize}

\subsection{Curvature-squared basis}

In the same reduced basis, ORBIT finds
\begin{align}
v_{R^2} &= (-2,1,1,0,0),\\
v_{R_{\mu\nu}^2} &= \left(-\frac12,\frac14,\frac12,-\frac12,\frac14\right).
\end{align}

\subsection{Independent quadratic redefinition vectors}

The independent quadratic higher-derivative image vectors are
\begin{align}
r_1 &= (-2,0,2,0,0),\\
r_2 &= (2,-1,0,-2,1),\\
r_3 &= (2,-2,0,0,0).
\end{align}

These come from the three independent quadratic redefinitions
\begin{align}
\delta h^{(1)}_{\mu\nu} &= \eta_{\mu\nu}\,\partial_\alpha\partial_\beta h^{\alpha\beta},\\
\delta h^{(2)}_{\mu\nu} &= \Box h_{\mu\nu},\\
\delta h^{(3)}_{\mu\nu} &= \eta_{\mu\nu}\,\Box h.
\end{align}

\subsection{The linear system}

We now solve
\begin{equation}
v_{2,4}
=
c_{R^2}v_{R^2}
c_{R_{\mu\nu}^2}v_{R_{\mu\nu}^2}
\alpha r_1+\beta r_2+\gamma r_3.
\end{equation}

The stored solution is
\begin{equation}
\boxed{
c_{R^2}=0,\qquad
c_{R_{\mu\nu}^2}=0,\qquad
\alpha=-\frac{3}{10M^2},\qquad
\beta=-\frac{3}{20M^2},\qquad
\gamma=\frac{3}{10M^2}.}
\end{equation}

Therefore
\begin{equation}
v_{2,4}
=
-\frac{3}{10M^2}r_1
-\frac{3}{20M^2}r_2
+\frac{3}{10M^2}r_3.
\end{equation}

The reconstruction can be checked directly:
\begin{align}
-\frac{3}{10}r_1
-\frac{3}{20}r_2
+\frac{3}{10}r_3
&=
\left(
\frac{6}{10},0,-\frac{6}{10},0,0
\right)
\nonumber\\
&\quad+
\left(
-\frac{3}{10},
\frac{3}{20},
0,
\frac{3}{10},
-\frac{3}{20}
\right)
\nonumber\\
&\quad+
\left(
\frac{6}{10},
-\frac{6}{10},
0,0,0
\right)
\nonumber\\
&=
\left(
\frac{9}{10},
-\frac{9}{20},
-\frac35,
\frac{3}{10},
-\frac{3}{20}
\right).
\end{align}
After restoring the overall factor $1/M^2$, this is exactly $v_{2,4}$.

Since the curvature coefficients are zero, there is no surviving $R^2$ or $R_{\mu\nu}^2$ component. The entire leftover lies in the quadratic redefinition image.

\subsection{Explicit removing field redefinition}

One convenient removing redefinition is therefore
\begin{equation}
\boxed{
\delta h_{\mu\nu}
=
-\frac{3}{10M^2}\eta_{\mu\nu}\partial_\alpha\partial_\beta h^{\alpha\beta}
-\frac{3}{20M^2}\Box h_{\mu\nu}
+\frac{3}{10M^2}\eta_{\mu\nu}\Box h.}
\end{equation}

Applying the induced quadratic variation of the Fierz--Pauli action removes the full $\{2,4\}$ remainder.

\subsection{Direct operator-level cancellation}

To make the last step completely explicit, introduce the five reduced-basis operators
\begin{align}
O_1 &:= \partial^\gamma\partial^\beta h^\alpha{}_\alpha\,
\partial_\delta\partial^\delta h_{\beta\gamma},\\
O_2 &:= \partial_\beta\partial^\beta h^\alpha{}_\alpha\,
\partial_\delta\partial^\delta h^\gamma{}_\gamma,\\
O_3 &:= \partial_\beta\partial_\alpha h_{\gamma\delta}\,
\partial^\delta\partial^\gamma h^{\alpha\beta},\\
O_4 &:= \partial_\delta\partial_\beta h_{\alpha\gamma}\,
\partial^\delta\partial^\gamma h^{\alpha\beta},\\
O_5 &:= \partial_\delta\partial_\gamma h_{\alpha\beta}\,
\partial^\delta\partial^\gamma h^{\alpha\beta}.
\end{align}
These are exactly the five reduced basis elements in which ORBIT stores the $\{2,4\}$ coordinates.

After substituting the matched branch, the leftover unsupported difference is
\begin{equation}
\Delta\mathcal L_{2,4}
=
\frac{1}{M^2}\left(
\frac{9}{10}O_1
-\frac{9}{20}O_2
-\frac35 O_3
+\frac{3}{10}O_4
-\frac{3}{20}O_5
\right).
\end{equation}

Now split the removing redefinition into three pieces,
\begin{align}
\delta h_{\mu\nu}^{(1)} &:= \eta_{\mu\nu}\partial_\alpha\partial_\beta h^{\alpha\beta},\\
\delta h_{\mu\nu}^{(2)} &:= \Box h_{\mu\nu},\\
\delta h_{\mu\nu}^{(3)} &:= \eta_{\mu\nu}\Box h,
\end{align}
so that
\[
\delta h_{\mu\nu}
=
-\frac{3}{10M^2}\delta h_{\mu\nu}^{(1)}
-\frac{3}{20M^2}\delta h_{\mu\nu}^{(2)}
+\frac{3}{10M^2}\delta h_{\mu\nu}^{(3)}.
\]

For the quadratic action,
\[
\delta\mathcal L_{\mathrm{FP}}
\simeq
E_{\mathrm{FP}}^{\mu\nu}\delta h_{\mu\nu},
\]
where $\simeq$ means equality modulo total derivatives. Reducing each image to the same $\{2,4\}$ basis gives
\begin{align}
E_{\mathrm{FP}}^{\mu\nu}\delta h_{\mu\nu}^{(1)}
&\simeq
-2O_1+2O_3,\\
E_{\mathrm{FP}}^{\mu\nu}\delta h_{\mu\nu}^{(2)}
&\simeq
2O_1-O_2-2O_4+O_5,\\
E_{\mathrm{FP}}^{\mu\nu}\delta h_{\mu\nu}^{(3)}
&\simeq
2O_1-2O_2.
\end{align}

Therefore
\begin{align}
\delta\mathcal L_{\mathrm{FP}}
&\simeq
-\frac{3}{10M^2}(-2O_1+2O_3)
-\frac{3}{20M^2}(2O_1-O_2-2O_4+O_5)
+\frac{3}{10M^2}(2O_1-2O_2)\\
&=
\frac{1}{M^2}\left(
\frac{9}{10}O_1
-\frac{9}{20}O_2
-\frac35 O_3
+\frac{3}{10}O_4
-\frac{3}{20}O_5
\right)\\
&=
\Delta\mathcal L_{2,4}.
\end{align}

Hence the transformed leftover is
\begin{equation}
\Delta\mathcal L_{2,4}-\delta\mathcal L_{\mathrm{FP}} \simeq 0.
\end{equation}

This is the direct ``on paper'' statement that the apparent curvature-squared remainder does not survive: every operator in the $\{2,4\}$ remainder is cancelled by the explicit quadratic higher-derivative field redefinition above.

\section{Final Conclusion}

The comparison is therefore complete.

\subsection*{Supported sectors}
After canonical normalization and the solved branch
\[
\boxed{
\Lmu=-\frac32,\qquad
\Lambda_{\mathrm{ref}}=0,\qquad
\kappa_{\mathrm{ref}}=\frac{4\sqrt3}{M},
\qquad
M^2=e^{3/2}\bar\mu^2,}
\]
all supported reduced sectors
\[
\{1,0\},\{2,0\},\{2,2\},\{3,0\},\{3,2\},\{4,0\},\{4,2\},\{5,0\},\{6,0\}
\]
match the sign-flipped Einstein--Hilbert reference exactly.

\subsection*{Quadratic four-derivative sector}
The only apparent leftover sector, $\{2,4\}$, has zero curvature-squared component and is entirely a quadratic higher-derivative field-redefinition image.

\subsection*{The completed statement}
Hence the canonically normalized Matchete-derived EFT is equivalent through mass dimension six to the sign-flipped Einstein--Hilbert action with vanishing cosmological constant, once one quotients by:
\begin{itemize}
\item integration by parts,
\item the interaction-sector Fierz--Pauli field-redefinition image already implemented by ORBIT,
\item the quadratic higher-derivative redefinitions needed to remove the $\{2,4\}$ remainder.
\end{itemize}

\end{document}
